% !TEX program = xelatex
% =====================================================================
% LLM Theory Seminar Week 5: Transformer Computational Theory
%
% Compile:
%   xelatex transformer_computational_theory_notes.tex
%   xelatex transformer_computational_theory_notes.tex
% or
%   latexmk -xelatex transformer_computational_theory_notes.tex
%
% Korean typesetting: kotex + XeLaTeX.
%
% Fixed layout reference / inspiration:
%   - Ben Jones, "lecture-notes" (Overleaf, CC BY 4.0)
%   - CMU-Notes/notes-template (structured university course notes)
% This file intentionally keeps the approved Week-1/2/3 seminar-note layout.
% =====================================================================

\documentclass[11pt,a4paper,oneside,openany]{memoir}

\usepackage{kotex}
\usepackage{amsmath,amssymb,amsthm,mathtools,bm}
\usepackage{booktabs,longtable,array,tabularx,makecell}
\newcolumntype{Y}{>{\raggedright\arraybackslash}X}
\usepackage{enumitem}
\usepackage[most]{tcolorbox}
\usepackage[dvipsnames]{xcolor}
\usepackage{xurl}
\usepackage[unicode]{hyperref}
\usepackage{cleveref}
\usepackage{needspace}

% ---------- page geometry ----------
\setlrmarginsandblock{27mm}{27mm}{*}
\setulmarginsandblock{24mm}{28mm}{*}
\checkandfixthelayout

\setlength{\parindent}{0pt}
\setlength{\parskip}{0.48em}
\setlength{\emergencystretch}{2em}
\setlength{\headheight}{15pt}
\linespread{1.055}\selectfont
\setlist[itemize]{topsep=0.28em,itemsep=0.12em,parsep=0pt,leftmargin=1.55em}
\setlist[enumerate]{topsep=0.28em,itemsep=0.16em,parsep=0pt,leftmargin=1.75em}

% ---------- restrained lecture-note palette ----------
\definecolor{NoteBlue}{HTML}{294E6B}
\definecolor{NoteBlueLight}{HTML}{F2F6F9}
\definecolor{NoteGray}{HTML}{5A6268}
\definecolor{NoteGrayLight}{HTML}{F6F6F4}
\definecolor{NoteWarm}{HTML}{7A6545}
\definecolor{NoteWarmLight}{HTML}{FAF7F0}
\definecolor{NoteRule}{HTML}{C8CDD2}

\hypersetup{
  colorlinks=true,
  linkcolor=NoteBlue,
  citecolor=NoteBlue,
  urlcolor=NoteBlue,
  pdftitle={Transformer Computational Theory},
  pdfauthor={LLM Theory Seminar}
}

% ---------- running heads ----------
\makepagestyle{seminar}
\makeoddhead{seminar}{\footnotesize LLM Theory Seminar}{ }{\footnotesize Transformer Computational Theory}
\makeoddfoot{seminar}{}{\footnotesize\thepage}{}
\makeheadrule{seminar}{\textwidth}{0.35pt}
\pagestyle{seminar}
\aliaspagestyle{chapter}{seminar}

% ---------- chapter / section hierarchy ----------
\makechapterstyle{seminarnotes}{
  \setlength{\beforechapskip}{8pt}
  \setlength{\midchapskip}{8pt}
  \setlength{\afterchapskip}{22pt}
  \renewcommand{\chapterheadstart}{\vspace*{\beforechapskip}}
  \renewcommand{\chapnamefont}{\normalfont\small\bfseries}
  \renewcommand{\chapnumfont}{\normalfont\bfseries\fontsize{34}{38}\selectfont\color{NoteBlue}}
  \renewcommand{\chaptitlefont}{\normalfont\huge\bfseries\color{black}}
  \renewcommand{\printchaptername}{}
  \renewcommand{\chapternamenum}{}
  \renewcommand{\printchapternum}{\chapnumfont\thechapter}
  \renewcommand{\afterchapternum}{\par\nobreak\color{black}\vskip\midchapskip}
  \renewcommand{\printchaptertitle}[1]{\chaptitlefont ##1}
  \renewcommand{\afterchaptertitle}{\par\nobreak\vskip\afterchapskip}
}
\chapterstyle{seminarnotes}
\setsecheadstyle{\Large\bfseries}
\setsubsecheadstyle{\large\bfseries}
\setsubsubsecheadstyle{\normalsize\bfseries}
\setbeforesecskip{-1.3\baselineskip plus -0.2\baselineskip minus -0.1\baselineskip}
\setaftersecskip{0.55\baselineskip}
\setbeforesubsecskip{-0.9\baselineskip plus -0.15\baselineskip minus -0.1\baselineskip}
\setaftersubsecskip{0.35\baselineskip}

% ---------- notation ----------
\newcommand{\R}{\mathbb{R}}
\newcommand{\N}{\mathbb{N}}
\newcommand{\E}{\mathbb{E}}
\newcommand{\X}{\mathbf{X}}
\newcommand{\Q}{\mathbf{Q}}
\newcommand{\K}{\mathbf{K}}
\newcommand{\V}{\mathbf{V}}
\newcommand{\A}{\mathbf{A}}
\newcommand{\W}{\mathbf{W}}
\newcommand{\one}{\mathbf{1}}
\newcommand{\softmax}{\operatorname{softmax}}
\newcommand{\hardmax}{\operatorname{hardmax}}
\newcommand{\ReLU}{\operatorname{ReLU}}
\newcommand{\Attn}{\operatorname{Attn}}
\newcommand{\FF}{\operatorname{FF}}
\newcommand{\rank}{\operatorname{rank}}
\newcommand{\diag}{\operatorname{diag}}
\newcommand{\norm}[1]{\left\lVert #1\right\rVert}
\newcommand{\abs}[1]{\left\lvert #1\right\rvert}
\newcommand{\inner}[2]{\left\langle #1,#2\right\rangle}
\newcommand{\dd}{\,\mathrm{d}}
\DeclareMathOperator*{\argmax}{arg\,max}
\DeclareMathOperator*{\argmin}{arg\,min}

% ---------- note environments ----------
\tcbset{
  seminarbox/.style={
    enhanced,breakable,sharp corners,boxrule=0pt,frame hidden,
    left=3.2mm,right=3mm,top=1.8mm,bottom=1.8mm,
    before={\par\Needspace{5\baselineskip}},before skip=0.8em,after skip=0.8em,
    fonttitle=\bfseries\small,coltitle=black,
    attach title to upper={\quad},
  }
}
\newtcolorbox{keybox}[1][]{
  seminarbox,colback=NoteBlueLight,
  borderline west={1.8pt}{0pt}{NoteBlue},
  title={핵심#1}
}
\newtcolorbox{paperbox}[1][]{
  seminarbox,colback=NoteGrayLight,
  borderline west={1.8pt}{0pt}{NoteGray},
  title={원 논문 기준#1}
}
\newtcolorbox{suppbox}[1][]{
  seminarbox,colback=NoteWarmLight,
  borderline west={1.8pt}{0pt}{NoteWarm},
  title={보충 유도 --- 논문 밖의 독자용 전개#1}
}
\newtcolorbox{warningbox}[1][]{
  seminarbox,colback=NoteGrayLight,
  borderline west={1.8pt}{0pt}{NoteGray},
  title={주의#1}
}
\newtcolorbox{presentbox}{
  seminarbox,colback=white,
  borderline west={2.2pt}{0pt}{black!70},
  fontupper=\footnotesize,
  top=1.2mm,bottom=1.2mm,before skip=0.5em,after skip=0.5em,
  title={발표에서 이렇게 말하기}
}
\newtcolorbox{presentboxcompact}{
  enhanced,sharp corners,boxrule=0pt,frame hidden,colback=white,
  borderline west={2.2pt}{0pt}{black!70},
  left=3.2mm,right=3mm,top=0.8mm,bottom=0.8mm,
  before={\par},before skip=0.35em,after skip=0.35em,
  fonttitle=\bfseries\small,coltitle=black,
  fontupper=\footnotesize,
  attach title to upper={\quad},
  title={발표에서 이렇게 말하기}
}

\theoremstyle{definition}
\newtheorem{definition}{정의}[chapter]
\newtheorem{proposition}{명제}[chapter]
\newtheorem{theorem}{정리}[chapter]
\newtheorem{lemma}{보조정리}[chapter]


% ---------- Week 5 complexity notation ----------
\newcommand{\bits}{\{0,1\}}
\newcommand{\poly}{\operatorname{poly}}
\newcommand{\TIME}{\mathsf{TIME}}
\newcommand{\SPACE}{\mathsf{SPACE}}
\newcommand{\AC}{\mathsf{AC}}
\newcommand{\TC}{\mathsf{TC}}
\newcommand{\NC}{\mathsf{NC}}
\newcommand{\Reg}{\mathsf{REG}}
\newcommand{\CFL}{\mathsf{CFL}}
\newcommand{\CSL}{\mathsf{CSL}}
\newcommand{\CoT}{\mathsf{CoT}}
\newcommand{\Lclass}{\mathsf{L}}
\newcommand{\NL}{\mathsf{NL}}
\newcommand{\Pclass}{\mathsf{P}}


\begin{document}

\begin{titlingpage}
\thispagestyle{empty}
\vspace*{1.2cm}
{\small\bfseries LLM Theory Seminar \hfill Week 5\par}
\vspace{1.4cm}
{\HUGE\bfseries Transformer Computational Theory\par}
\vspace{0.55cm}
{\large formal languages, circuit complexity, Turing completeness, chain-of-thought computation\par}
\vspace{0.9cm}
\rule{\textwidth}{0.7pt}
\vspace{1.1cm}

\begin{minipage}{0.86\textwidth}
\small
이번 주의 질문은 ``Transformer가 어떤 연속함수를 근사할 수 있는가?''에서 한 단계 더 나아간다.
\textbf{입력 길이가 커질 때, 제한된 depth와 precision을 가진 Transformer가 실제로 어떤 계산 문제를 풀 수 있는가?}
이 질문에는 서로 다른 답이 공존한다. 이상화된 arbitrary-precision hard-attention Transformer는 Turing complete일 수 있지만,
현실적인 log-precision의 고정-depth Transformer 한 번의 forward pass는 uniform $\TC^0$로 upper bound된다.
그리고 autoregressive intermediate generation을 허용하면 생성 step 수 자체가 새로운 순차 계산 자원이 되어,
선형 길이의 scratchpad는 모든 regular language를, 다항 길이의 scratchpad는 적절한 가정 아래 정확히 $\Pclass$ 수준의 계산을 지원한다.
이번 노트의 목적은 이 결과들이 서로 모순이 아니라 \textbf{서로 다른 resource model}에 대한 정리라는 점을 수식으로 분리하는 것이다.
\end{minipage}

\vfill
{\small
\textbf{Primary readings}\\[0.25em]
Jorge Pérez, Pablo Barceló, Javier Marinkovic,
\textbf{Attention is Turing Complete}, JMLR 2021.\\
William Merrill, Ashish Sabharwal,
\textbf{The Parallelism Tradeoff: Limitations of Log-Precision Transformers}, TACL 2023.\\
William Merrill, Ashish Sabharwal,
\textbf{The Expressive Power of Transformers with Chain of Thought}, ICLR 2024.\\
Michael Hahn,
\textbf{Theoretical Limitations of Self-Attention in Neural Sequence Models}, TACL 2020.
\par}
\vspace{0.8cm}
{\small September 2026\par}
\end{titlingpage}

\tableofcontents*
\clearpage

% =====================================================================
\chapter{먼저 맞춰 둘 기호와 전체 그림}
% =====================================================================

\section{이번 주에 섞이면 안 되는 네 개의 축}

Transformer computational theory에서 가장 흔한 혼동은 서로 다른 resource를 한 문장에 섞는 것이다.
같은 ``Transformer''라는 단어라도 theorem마다 다음 네 가지가 다르다.

\begin{center}
\small
\begin{tabularx}{\textwidth}{>{\bfseries\raggedright\arraybackslash}p{0.18\textwidth}Y Y}
\toprule
축 & 대표 선택지 & 왜 중요한가 \\
\midrule
Attention & soft / hard / saturated & exact addressing, averaging, circuit simulation 가능성이 달라진다. \\
Precision & $O(1)$ / $O(\log n)$ / arbitrary & 한 scalar가 저장할 수 있는 정보량이 달라진다. \\
Sequential steps & one pass / $t(n)$ decoding steps & fixed-depth parallel 계산인지, 반복적인 순차 계산인지가 달라진다. \\
Architecture scaling & depth/width fixed vs.\ $n$에 따라 증가 & ``하나의 모델''인지 input length마다 더 큰 회로를 허용하는지 결정한다. \\
\bottomrule
\end{tabularx}
\end{center}

\begin{keybox}
``Transformer는 Turing complete다''와 ``Transformer는 $\TC^0$를 넘지 못한다''는 문장은 동시에 참일 수 있다.
앞 문장은 arbitrary precision, hard attention, unbounded decoding을 허용하는 계산 모델에 관한 결과이고,
뒤 문장은 log precision과 고정된 forward-pass depth를 가진 모델에 관한 upper bound다.
\end{keybox}

\section{논문별 기호 대응표}

\begin{center}
\small
\begin{tabularx}{\textwidth}{>{\bfseries\raggedright\arraybackslash}p{0.17\textwidth}Y Y Y}
\toprule
이 노트 & Pérez et al. & Merrill--Sabharwal & 의미 \\
\midrule
$n$ & $|w|$ & $n$ & input length \\
$L$ & encoder/decoder layer 수 & $d$ & Transformer depth \\
$H$ & attention head 수 & $h$ & number of heads \\
$p(n)$ & rational precision & $p$ & arithmetic precision in bits \\
$t(n)$ & output length / TM steps & $t(n)$ & intermediate decoding steps \\
$\Sigma$ & $\Sigma$ & $\Sigma$ & finite alphabet \\
$\mathcal L$ & $L$ & $L$ & formal language \\
$M$ & $M=(Q,\Sigma,\delta,q_{\mathrm{init}},F)$ & Turing machine $M$ & simulated machine \\
\bottomrule
\end{tabularx}
\end{center}

\section{결론을 먼저 보는 지도}

이번 주 결과를 resource 관점에서 한 줄씩 적으면 다음과 같다.

\begin{align*}
\text{hard/saturated, fixed depth}
&\leadsto \AC^0\text{ 또는 }\TC^0\text{ upper bound},\\
\text{log precision, one forward pass}
&\leadsto \text{uniform }\TC^0,\\
\text{arbitrary precision + hard attention}
&\leadsto \text{Turing completeness with unbounded decoding},\\
t(n)\text{ autoregressive steps}
&\leadsto t(n)\text{이 순차 계산 resource가 됨}.
\end{align*}

따라서 이번 주의 핵심 질문은 architecture 이름이 아니라
\[
\boxed{\text{무엇을 고정하고, 무엇을 input length와 함께 증가시키는가?}}
\]
이다.

\begin{presentbox}
``이번 주에는 Transformer의 계산력을 하나의 숫자로 말하지 않겠습니다. Attention type, numerical precision, model depth, 그리고 autoregressive step 수를 따로 잡아야 theorem들이 서로 맞아떨어집니다. 특히 one-pass Transformer와 CoT를 쓰는 decoder는 계산 모델 자체가 다릅니다.''
\end{presentbox}

% =====================================================================
\chapter{Formal language와 circuit complexity의 최소 배경}
% =====================================================================

\section{Language recognition을 계산 문제로 보기}

\begin{definition}[Formal language]
유한 alphabet $\Sigma$에 대해
\[
\Sigma^*=\bigcup_{n\ge 0}\Sigma^n
\]
이고, formal language는 임의의 부분집합
\[
\mathcal L\subseteq \Sigma^*
\]
이다.
\end{definition}

모델이 문자열 $x\in\Sigma^*$를 받아
\[
f(x)\in\{0,1\}
\]
을 출력하고
\[
f(x)=1 \iff x\in\mathcal L
\]
이면 $f$가 $\mathcal L$을 recognize한다고 한다.

이 정의는 language modeling probability와 다르다. 여기서는 특정 distribution 위 평균 loss가 아니라
\textbf{모든 길이의 모든 입력에 대해 membership을 정확히 계산할 수 있는지}를 본다.

\section{Finite automaton, regular language, 그리고 sequentiality}

DFA는
\[
A=(Q,\Sigma,\delta,q_0,F)
\]
로 쓰며 입력
\[
x=x_1x_2\cdots x_n
\]
에 대해
\[
q_i=\delta(q_{i-1},x_i),\qquad i=1,\dots,n
\]
을 반복한다.

최종적으로
\[
x\in\mathcal L(A)\iff q_n\in F.
\]

즉 DFA 자체는 아주 작은 memory만 필요하지만 계산 graph는 길이 $n$의 recurrence다.
Transformer computational theory에서 이것이 중요한 이유는 ``작은 memory''와 ``작은 parallel depth''가 같은 개념이 아니기 때문이다.

\begin{suppbox}
Transition을 함수 합성으로 쓰면
\[
q_n=\delta_{x_n}\circ\delta_{x_{n-1}}\circ\cdots\circ\delta_{x_1}(q_0),
\qquad
\delta_a(q):=\delta(q,a).
\]
순차적으로는 $n$번 적용하지만, 함수 합성이 결합법칙을 가지므로 parallel prefix 방식으로 tree를 만들면
$O(\log n)$ depth로도 계산 가능하다. 이 관점이 이후 automata shortcut과 $\NC^1$ 연결의 직관이다.
\end{suppbox}

\section{Circuit family: input length마다 회로 하나}

Boolean circuit $C_n$은 길이 $n$인 bit string을 입력받는다. 모든 길이를 처리하려면
\[
\{C_n\}_{n\ge 1}
\]
이라는 circuit family를 생각한다.

두 resource가 핵심이다.
\begin{itemize}
  \item \textbf{size}: gate/wire의 총 개수,
  \item \textbf{depth}: input에서 output까지 가장 긴 dependency chain.
\end{itemize}

Transformer layer 수가 input length에 무관한 상수라면, forward pass를 회로로 펼쳤을 때 자연스럽게 constant-depth computation과 비교하게 된다.

\section{\texorpdfstring{$\AC^0$, $\TC^0$, $\NC^1$}{AC0, TC0, NC1}}

\begin{definition}[$\AC^0$]
Polynomial-size, constant-depth Boolean circuit family이며 AND, OR, NOT gate가 허용된다.
AND/OR의 fan-in은 unbounded일 수 있다.
\end{definition}

\begin{definition}[$\TC^0$]
$\AC^0$에 unbounded-fan-in threshold/majority gate를 허용한 클래스다. 대표 threshold gate는
\[
\operatorname{THR}(x_1,\dots,x_m)
=
\one\!\left[\sum_{i=1}^{m}w_i x_i\ge \theta\right].
\]
\end{definition}

표준 uniform 버전을 염두에 두면 다음 containment를 기억하면 충분하다.
\[
\AC^0\subsetneq \TC^0\subseteq \NC^1\subseteq \Lclass\subseteq \NL\subseteq \Pclass.
\]

첫 번째 strictness는 parity/majority에 대한 고전적 lower bound에서 오지만,
$\TC^0$와 $\NC^1$ 사이의 strictness 같은 더 뒤의 분리는 알려져 있지 않다.
따라서 논문은 보통
\[
\TC^0\neq \NC^1,\qquad \Lclass\neq\Pclass
\]
같은 표준 complexity conjecture를 조건으로 ``못 푼다''는 결론을 말한다.

\section{\texorpdfstring{왜 $\TC^0$가 생각보다 약하지 않은가}{왜 TC0가 생각보다 약하지 않은가}}

$\TC^0$는 constant depth라고 해서 단순한 함수만 포함하는 것은 아니다.
threshold gate 덕분에 다수결, 정수 덧셈, 곱셈과 같은 상당히 풍부한 arithmetic이 매우 parallel하게 구현될 수 있다.

따라서
\[
\text{``Transformer}\in\TC^0\text{''}
\]
라는 upper bound는 ``덧셈도 못한다''가 아니라
즉 이 upper bound의 해석은
\begin{quote}
``본질적으로 긴 sequential dependency가 필요한 일부 문제에서는
한 번의 고정-depth pass가 제한을 받는다.''
\end{quote}
에 가깝다.

\begin{warningbox}
Complexity class inclusion은 실제 LLM accuracy에 대한 직접 예측이 아니다. Theorem의 모델은 precision, activation, masking, uniformity를 엄밀히 고정한 계산 모델이다. 실제 finite-length benchmark에서 특정 task를 근사적으로 맞히는 것과 모든 $n$에 대해 exact recognition하는 것은 다른 주장이다.
\end{warningbox}

\begin{presentbox}
``Formal language를 쓰는 이유는 자연어가 정규언어라는 주장을 하려는 게 아닙니다. 입력 길이가 커질 때 모델이 parity, bracket matching, automaton simulation 같은 계산을 정확히 수행할 수 있는지를 clean하게 측정하기 위해서입니다. 그리고 circuit complexity는 Transformer의 강한 parallelism을 depth라는 자원으로 번역해 줍니다.''
\end{presentbox}

% =====================================================================
\chapter{고정-depth self-attention의 초기 한계}
% =====================================================================

\section{Hahn의 질문: recurrence 없이 hierarchy를 처리할 수 있는가?}

Hahn (2020)은 고정된 layer/head 수의 self-attention 모델이 input length가 증가할 때
formal language를 어떻게 처리하는지 분석했다. 대표 예는 다음 두 계산이다.

\begin{itemize}
  \item \textbf{Parity}: $x\in\bits^n$에 대해
  \[
  \operatorname{PARITY}(x)=\left(\sum_{i=1}^n x_i\right)\bmod 2.
  \]
  \item \textbf{Dyck language}: 괄호 prefix마다 열린 괄호 수가 음수가 되지 않고 마지막에 0이 되는지 검사한다.
\end{itemize}

Parity는 finite-state computation이지만 모든 bit의 global interaction이 필요하고,
Dyck는 unbounded hierarchical memory를 요구한다. Hahn의 결과는 fixed-size self-attention이 이런 구조를 길이에 무관하게 exact하게 처리하는 데 강한 제약이 있음을 보였다.

\begin{paperbox}[--- Hahn (2020)]
Soft attention과 hard attention 모두에서, model의 layer 수나 head 수를 input length에 따라 키우지 않는 한 periodic finite-state language와 hierarchical structure를 일반적으로 처리할 수 없다는 한계가 제시된다. 이 결과는 이후 circuit-complexity upper bound로 더 구조적으로 정리된다.
\end{paperbox}

\section{Hard attention을 회로로 보면 더 깔끔해진다}

Hao, Angluin, Frank (2022)는 hard-attention Transformer를 세 종류로 나눈다.

\begin{itemize}
  \item \textbf{UHAT}: unique hard attention,
  \item \textbf{GUHAT}: UHAT의 generalized version,
  \item \textbf{AHAT}: maximum score 위치들을 평균하는 averaging hard attention.
\end{itemize}

핵심 upper bound는
\[
\boxed{\mathrm{UHAT},\mathrm{GUHAT}\subseteq \AC^0}
\]
이다.

$\AC^0$는 parity를 recognize할 수 없으므로 즉시
\[
\mathrm{PARITY}\notin \mathcal L(\mathrm{GUHAT})
\]
같은 결론이 따른다.

반면 AHAT는 tie된 여러 위치를 평균할 수 있어 더 강해지고, 논문에서는 MAJORITY와 DYCK-1 같은 non-$\AC^0$ language를 recognize하는 construction을 보인다.

\section{``attention이 global이니까 sequential problem도 쉽다''가 왜 틀릴 수 있는가}

Self-attention은 한 층에서 모든 token pair를 연결한다.

\[
a_i
=
\sum_{j=1}^{n}\alpha_{ij}v_j.
\]

하지만 dependency graph가 dense하다는 사실과 arbitrary sequential algorithm을 구현할 수 있다는 사실은 다르다.
한 층은 모든 값을 \textbf{볼 수} 있지만, 여러 단계의 state transition을 차례로 \textbf{적용}해야 하는 계산은 별도의 depth 또는 recurrence가 필요할 수 있다.

예를 들어 DFA transition을
\[
q_n
=
\delta_{x_n}\circ\cdots\circ\delta_{x_1}(q_0)
\]
로 쓰면, 모든 $x_i$를 한 번에 읽는 것과 이 ordered composition을 계산하는 것은 다른 문제다.

\begin{keybox}
Attention의 ``global receptive field''는 communication distance를 줄여 주지만, computational depth를 자동으로 늘려 주지 않는다.
이번 주 이론에서 가장 중요한 구분 중 하나다.
\end{keybox}

\begin{presentbox}
``한 층의 attention은 모든 토큰을 볼 수 있지만, 모든 알고리즘을 한 단계에 수행할 수 있다는 뜻은 아닙니다. Circuit 관점에서는 communication은 global이어도 layer 수가 고정되어 있으면 computational depth는 여전히 상수입니다. 그래서 parity나 automaton composition 같은 문제가 중요한 test case가 됩니다.''
\end{presentbox}

% =====================================================================
\chapter{\texorpdfstring{Log-precision Transformer와 $\TC^0$ upper bound}{Log-precision Transformer와 TC0 upper bound}}
% =====================================================================

\section{왜 precision을 명시해야 하는가?}

실수 하나에 arbitrary precision을 허용하면 이론적으로 무한히 많은 정보를 한 scalar에 숨길 수 있다.
따라서 실제 digital computation에 가까운 모델을 만들려면 input length $n$에 따라 사용할 수 있는 bit 수를 제한해야 한다.

Merrill--Sabharwal (2023)의 핵심 setting은
\[
p(n)=O(\log n)
\]
bit precision이다.

이 정도 precision이면 값의 표현 범위와 resolution이 polynomial scale에 머물면서도,
길이 $n$의 sequence 위 uniform averaging 같은 attention operation을 표현할 수 있다.

\section{Main theorem}

\begin{theorem}[Merrill--Sabharwal 2023, informal]
Arithmetic precision이 $O(\log n)$이고 feed-forward subnetwork가 입력 크기에 선형인 space에서 계산 가능한 Transformer classifier는
constant-depth logspace-uniform threshold circuit로 simulate할 수 있다. 따라서
\[
\boxed{\text{log-precision fixed-depth Transformer}\subseteq \text{uniform }\TC^0.}
\]
\end{theorem}

\begin{paperbox}
이 theorem에서 ``Transformer''는 input length가 커져도 layer 수, head 수, hidden architecture가 고정된 family를 의미한다. 핵심은 한 번의 forward pass를 constant-depth threshold circuit로 compile할 수 있다는 것이다.
\end{paperbox}

\section{왜 circuit depth가 constant인가? --- proof skeleton}

Transformer layer를 대략
\[
H^{(\ell+1)}
=
\operatorname{FFN}^{(\ell)}
\left(
H^{(\ell)}+
\operatorname{Attn}^{(\ell)}(H^{(\ell)})
\right)
\]
라고 하자.

전체 layer 수가 $L=O(1)$이면 다음을 보이면 충분하다.

\begin{enumerate}
  \item embedding과 positional feature가 $\TC^0$에서 계산 가능,
  \item attention score 계산이 $\TC^0$에서 가능,
  \item $n$개의 값의 합/정규화가 $\TC^0$에서 가능,
  \item FFN arithmetic이 $\TC^0$에서 가능,
  \item residual/normalization도 $\TC^0$에서 가능.
\end{enumerate}

각 block이 constant-depth threshold circuit이고 layer 수 $L$도 constant이면 합성한 depth는
\[
D_{\mathrm{total}}
\le
L\,D_{\mathrm{layer}}+D_{\mathrm{embed}}+D_{\mathrm{out}}
=O(1).
\]

또 각 위치 $i=1,\dots,n$의 계산을 parallel하게 복제하므로 전체 gate 수는 polynomial size에 머문다.

\begin{suppbox}
왜 $O(\log n)$ precision이 polynomial-size circuit와 잘 맞는지 보자. $p=c\log n$ bit이면 한 scalar는 최대
\[
2^p=n^c
\]
개의 bit pattern 수준으로 표현된다. 따라서 arithmetic operand의 bit-length는 input length에 대해 logarithmic한 범위에 머문다.

이 때문에 polynomial-size circuit simulation을 유지할 수 있다. 반대로 arbitrary precision을 허용하면 scalar 하나가 input length에 비해 과도하게 많은 정보를 저장할 수 있어 이 argument가 깨진다.
\end{suppbox}

\section{Self-attention의 합이 왜 ``순차 합산''을 요구하지 않는가?}

attention output은
\[
a_i=\frac{\sum_{j=1}^{n}s_{ij}v_j}{\sum_{j=1}^{n}s_{ij}}
\]
형태다. 프로그램으로 쓰면 $j=1$부터 $n$까지 loop를 도는 것처럼 보인다.
하지만 circuit model에서는 $n$개 항을 parallel tree 또는 threshold arithmetic circuit로 합산할 수 있다.

따라서 wall-clock sequential code와 circuit depth를 혼동하면 안 된다.
Transformer의 parallelism이 바로 이 지점에서 complexity upper bound에 반영된다.

\section{Upper bound의 실제 의미}

$\TC^0$ upper bound로부터 바로 ``특정 문제는 불가능''이라고 말하려면 complexity separation이 필요하다.
예를 들어 paper는 표준 가정 아래 다음 식의 해석을 강조한다.

\[
\text{one-pass Transformer}\in\TC^0
\quad\text{vs.}\quad
\text{problem complete for a larger class}.
\]

예를 들어 arbitrary finite-state machine simulation은 $\NC^1$ 수준의 sequential composition과 연결되고,
directed graph connectivity는 $\NL$-complete, 일부 algebraic reasoning problem은 $\Pclass$-complete하다.
따라서 적절한 class non-collapse assumption 아래 fixed-depth one-pass Transformer의 한계가 된다.

\begin{warningbox}
``$\TC^0$ upper bound이므로 모든 regular language를 못 푼다''라고 무조건 쓰면 안 된다. 많은 regular language는 $\AC^0$나 $\TC^0$ 안에 있다. 정확한 주장은, 표준 separation 가정 아래 \textbf{모든} regular language를 한 architecture class가 처리할 수는 없다는 식으로 해야 한다.
\end{warningbox}

\begin{presentbox}
``Log-precision theorem의 요지는 attention 연산 자체가 약하다는 게 아닙니다. 한 forward pass가 엄청나게 parallel하기 때문에, 그 계산 전체가 constant-depth threshold circuit로 compile됩니다. 이 parallelism이 장점인 동시에 sequential depth를 제한하는 원인이 됩니다.''
\end{presentbox}

% =====================================================================
\chapter{그런데 왜 Transformer가 Turing complete인가?}
% =====================================================================

\section{Turing completeness의 정의부터 분리한다}

\begin{definition}[Turing completeness, language-recognition view]
계산 모델의 class $\mathcal N$이 모든 Turing machine이 결정할 수 있는 language를 simulate할 수 있으면 $\mathcal N$을 Turing complete라고 한다.
\end{definition}

중요한 것은 Turing completeness가
\[
\text{``고정된 한 번의 forward pass로 모든 계산을 끝낸다''}
\]
는 뜻이 아니라는 점이다. Turing machine도 problem size에 따라 임의로 많은 step을 사용한다.

\section{Pérez--Barceló--Marinkovic의 main theorem}

\begin{theorem}[Pérez et al. 2021, Theorem 6]
Positional encoding을 가진 hard-attention Transformer class는 Turing complete하다. 논문의 construction에서는 positional embedding의 nonconstant coordinate로
\[
n,\qquad \frac1n,\qquad \frac1{n^2}
\]
만 사용하며, single encoder layer와 three decoder layers로도 충분하다.
\end{theorem}

논문은 더 강하게 recursively enumerable language까지 recognize할 수 있는 construction을 논의한다.

\begin{paperbox}[--- theorem의 중요한 가정]
증명은 rational-valued \textbf{arbitrary precision} internal representation과 \textbf{hard attention}을 사용한다. 저자들도 fixed precision이면 같은 Turing-completeness conclusion이 성립하지 않는다고 명시한다.
\end{paperbox}

\section{TM simulation의 상태를 무엇으로 저장하는가?}

한 tape Turing machine의 시각 $t$ configuration을 개념적으로
\[
C_t=(q_t,h_t,\tau_t)
\]
로 쓰자.

\begin{itemize}
  \item $q_t$: finite control state,
  \item $h_t\in\mathbb Z$: head position,
  \item $\tau_t:\mathbb Z\to\Sigma$: tape content.
\end{itemize}

transition function은
\[
\delta(q_t,\tau_t(h_t))
=
(q_{t+1},s_t^{\mathrm{write}},d_t),
\qquad d_t\in\{-1,+1\}
\]
이고
\[
h_{t+1}=h_t+d_t.
\]

문제는 tape 전체 $\tau_t$를 hidden vector 하나에 넣는 것이 아니다.
Decoder가 지금까지 생성한 history를 memory처럼 사용하고, hard attention으로 필요한 위치를 retrieve한다.

\section{가장 최근 write를 찾으면 tape를 읽을 수 있다}

현재 head가 cell $j=h_t$를 읽는다고 하자. 그 cell에 마지막으로 write한 시점을
\[
r_t(j)
:=
\max\{s<t: h_s=j\text{ and step }s\text{ writes to }j\}
\]
로 두자.

그러면 현재 tape symbol은
\[
\tau_t(j)=
\begin{cases}
\text{symbol written at step }r_t(j), & r_t(j)\text{가 존재},\\
\text{initial symbol at }j, & \text{otherwise}.
\end{cases}
\]

즉 decoder history 안에서
\[
\boxed{\text{``현재 cell과 같은 address를 가진 가장 최근 event''}}
\]
만 찾으면 된다. Pérez et al.의 hard-attention construction은 positional rational code를 이용해 이런 addressing을 구현한다.

\begin{suppbox}
이 construction을 database 관점으로 보면 이해가 쉽다. 각 decoding token은
\[
(\text{time},\text{head address},\text{write symbol},\text{state})
\]
를 기록하는 log entry다. 현재 step은 attention으로 과거 log 중 필요한 address의 최신 record를 읽고, FFN이 finite transition $\delta$를 적용해 다음 log entry를 쓴다. Autoregressive generation이 곧 Turing-machine time axis가 된다.
\end{suppbox}

\section{Positional encoding이 왜 핵심인가?}

Pérez et al.은 positional encoding이 없는 Transformer가 강한 proportion invariance를 가진다는 것을 먼저 보인다.
두 입력 $u,w$가 각 symbol의 비율을 같게 가지면 특정 설정에서 Transformer output이 같아지는 제약이 생긴다.

따라서 순서와 정확한 address를 다루려면 position 정보를 제공해야 한다.
Construction에서는 $n,1/n,1/n^2$ 형태를 사용해 decoder step과 tape address를 분리하고 비교할 수 있게 한다.

\section{Fixed precision이면 왜 같은 증명이 무너지는가?}

Arbitrary precision에서는 position index $n$이 아무리 커져도 서로 다른 rational code를 유지할 수 있다.
하지만 fixed precision에서는 positional value가 취할 수 있는 값의 집합이 finite다.

충분히 긴 sequence에서는 pigeonhole principle에 의해 서로 다른 position들이 같은 code로 충돌한다.
따라서 unbounded tape address를 exact하게 encode하는 Pérez-style proof가 유지되지 않는다.

\begin{keybox}
Turing completeness theorem의 ``infinite computation''은 hidden dimension이 무한해서가 아니다. 고정된 architecture가 arbitrary number of autoregressive steps와 arbitrary-precision numerical state를 사용해 Turing machine trajectory를 계속 이어 갈 수 있다는 뜻이다.
\end{keybox}

\begin{presentbox}
``Turing-complete proof에서 Transformer는 한 번에 Turing machine 전체를 계산하지 않습니다. Decoder가 한 token을 생성할 때마다 machine의 한 step을 진행하고, 이미 생성된 prefix를 외부 memory처럼 읽습니다. 그래서 이 결과를 one-pass $\TC^0$ upper bound와 비교하려면 decoding time과 precision을 반드시 같이 봐야 합니다.''
\end{presentbox}

% =====================================================================
\chapter{\texorpdfstring{모순 해소: one-pass $\TC^0$와 Turing completeness}{모순 해소: one-pass TC0와 Turing completeness}}
% =====================================================================

\section{두 theorem의 assumptions를 같은 표에 놓기}

\begin{center}
\small
\begin{tabularx}{\textwidth}{>{\bfseries\raggedright\arraybackslash}p{0.22\textwidth}Y Y}
\toprule
항목 & Log-precision upper bound & Pérez Turing completeness \\
\midrule
Precision & $O(\log n)$ bits & arbitrary-precision rationals \\
Attention & practical-style arithmetic / bounded precision & hard attention essential to proof \\
Sequential compute & one fixed-depth forward pass & unbounded autoregressive decoder steps \\
Position information & bounded representation & exact rational positional codes \\
Conclusion & uniform $\TC^0$ upper bound & arbitrary TM simulation 가능 \\
\bottomrule
\end{tabularx}
\end{center}

둘은 서로 다른 계산 모델이므로 논리적 충돌이 없다.

\section{계산 resource를 식으로 분해하면 더 명확하다}

Transformer decoder의 전체 계산량을 아주 거칠게
\[
\text{Total sequential depth}
\approx
\underbrace{L}_{\text{layers per pass}}
\times
\underbrace{T}_{\text{autoregressive passes}}
\]
로 생각할 수 있다.

one-pass theorem에서는
\[
L=O(1),\qquad T=O(1).
\]

Turing-machine simulation에서는 architecture depth $L$은 고정이어도
\[
T=T(n)
\]
가 input-dependent하게 커질 수 있다.

따라서 모델 parameter 수는 같아도 실행 시간이 계산력을 증가시킨다.

\section{Precision도 독립적인 resource다}

또 하나의 축은 bit precision이다.

\[
\text{information per scalar}\sim p(n)\text{ bits}.
\]

$O(\log n)$ precision은 polynomial range/addressing에 충분하지만 unbounded exact position encoding을 scalar 하나에 계속 넣을 수는 없다.
Arbitrary precision은 이 restriction을 없앤다.

따라서 computational power를 비교할 때 최소한
\[
(L,H,d,p(n),T(n))
\]
를 함께 적는 습관이 필요하다.

\begin{warningbox}
Turing completeness는 실용적 학습 가능성, numerical robustness, finite precision implementation을 보장하지 않는다. 반대로 $\TC^0$ upper bound도 finite benchmark에서 LLM이 heuristic이나 approximate solution으로 높은 accuracy를 낼 수 없다는 뜻이 아니다.
\end{warningbox}

\begin{presentbox}
``겉으로는 한쪽 논문이 Turing complete, 다른 논문이 $\TC^0$라고 해서 충돌해 보입니다. 하지만 계산 모델을 $(\text{precision},\text{depth},\text{decoding steps})$로 쓰면 바로 풀립니다. Turing-complete construction은 sequential time과 arbitrary precision을 자원으로 쓰고, $\TC^0$ theorem은 그 자원을 제한한 one-pass 모델입니다.''
\end{presentbox}

% =====================================================================
\chapter{Chain of Thought를 계산 resource로 보기}
% =====================================================================

\section{Intermediate token 하나는 forward pass 하나다}

Decoder-only Transformer를
\[
f:\Sigma^*\to\Sigma
\]
라고 하자. 한 번의 decoding으로
\[
y_1=f(x)
\]
를 만들고 다시 prefix에 붙이면
\[
y_2=f(xy_1),\qquad
\dots,\qquad y_t=f(xy_1\cdots y_{t-1}).
\]

즉 CoT/scratchpad 길이 $t$는 단순히 ``답을 길게 쓰는 것''이 아니라
\[
\boxed{t\text{번의 serial forward passes}}
\]
를 허용하는 resource다.

\begin{definition}[$\CoT(t(n))$]
입력 길이 $n$에 대해 최대 $t(n)$개의 intermediate decoding step을 사용하는 어떤 Transformer decoder가 recognize할 수 있는 language들의 class를 $\CoT(t(n))$로 쓴다.
\end{definition}

\section{Merrill--Sabharwal의 central sandwich}

ICLR 2024 결과를 핵심 형태로 쓰면, polylog factor를 숨겼을 때
\[
\boxed{
\TIME(t(n))
\subseteq
\CoT(t(n))
\subseteq
\SPACE(t(n)+\log n)
\subseteq
\widetilde{\TIME}(t(n)^2+n^2)
}
\]
이다.

여기서 왼쪽 lower bound construction은 strict causal saturated attention과 projected pre-norm을 사용한다.
Upper bound는 그 특수 construction에 의존하지 않는다.

\section{왼쪽 inclusion: TM step 하나를 token 하나로}

\[
\TIME(t(n))\subseteq\CoT(t(n))
\]
의 직관은 Pérez construction과 비슷하지만, log-precision decoder-only setting에 맞게 더 정교하게 만든 것이다.

Turing machine의 transition record를
\[
\Delta_i
=
(q_i,\text{symbols read},\text{symbols written},\text{head moves})
\]
라고 하자.

Construction은 induction으로
\[
\text{decoding step }i
\quad\mapsto\quad
\Delta_i
\]
를 생성한다.

현재 tape cell의 값을 알아내려면 가장 최근에 같은 cell에 write한 token을 찾아야 한다.
이를 위해 논문은 layer-norm hash를 사용한다.

\section{Layer-norm hash를 한 줄씩 유도}

논문의 핵심 device를 직접 계산해 보자.

\[
v_x:=\langle x,1,-x,-1\rangle.
\]
평균은 0이고
\[
\norm{v_x}^2
=x^2+1+x^2+1
=2x^2+2.
\]
따라서 normalized vector를
\[
\phi_x
:=
\frac{1}{\sqrt{2x^2+2}}
\langle x,1,-x,-1\rangle
\]
로 둔다.

이제 scale된 pair $(x/i,1/i)$로 같은 vector를 만들면
\begin{align*}
\phi(x/i,1/i)
&=
\frac{\langle x/i,1/i,-x/i,-1/i\rangle}
{\sqrt{2x^2/i^2+2/i^2}}\\
&=
\frac{(1/i)\langle x,1,-x,-1\rangle}
{(1/i)\sqrt{2x^2+2}}\\
&=
\phi_x.
\end{align*}

즉 layer normalization이 공통 scale $1/i$를 지워 주어 position-dependent denominator 속에서 discrete address $x$를 복원한다.

두 hash의 inner product는
\begin{align*}
\inner{\phi_q}{\phi_k}
&=
\frac{qk+1+qk+1}
{\sqrt{(2q^2+2)(2k^2+2)}}\\
&=
\frac{qk+1}
{\sqrt{(q^2+1)(k^2+1)}}.
\end{align*}

언제 이 값이 최대 $1$이 되는가?
\begin{align*}
1
&=
\frac{qk+1}{\sqrt{(q^2+1)(k^2+1)}}\\
(qk+1)^2
&=(q^2+1)(k^2+1)\\
q^2k^2+2qk+1
&=q^2k^2+q^2+k^2+1\\
0&=q^2-2qk+k^2\\
0&=(q-k)^2.
\end{align*}

따라서
\[
\boxed{\inner{\phi_q}{\phi_k}=1\iff q=k.}
\]

Attention score로 이 inner product를 쓰면 exact equality lookup을 구현할 수 있다.

\begin{paperbox}
이 layer-norm hash는 ICLR 2024 paper의 proof에서 핵심 construction이다. Projected pre-norm은 hidden state의 필요한 coordinate subset만 normalize해 이런 hash를 여러 용도로 만들 수 있게 한다.
\end{paperbox}

\section{Regular language는 왜 linear CoT면 충분한가?}

DFA state를 $q_i$라고 하면
\[
q_i=\delta(q_{i-1},x_i).
\]
Input을 다 읽은 뒤 intermediate token에서
\[
q_0,q_1,\dots,q_n
\]
을 차례로 생성하면 된다.

각 decoding step에서는
\begin{enumerate}
  \item 이전 state $q_{i-1}$를 current prefix에서 읽고,
  \item layer-norm hash로 $x_i$에 해당하는 input position을 찾고,
  \item FFN이 finite transition table $\delta$를 계산한다.
\end{enumerate}

따라서 모든 regular language $\mathcal L$에 대해 적절한 projected-pre-norm decoder가
\[
|x|+1
\]
decoding steps로 membership을 판정할 수 있다.

\section{\texorpdfstring{Polynomial CoT가 왜 $\Pclass$까지 가는가?}{Polynomial CoT가 왜 P까지 가는가?}}

Lower bound 쪽에서는 polynomial-time Turing machine의 각 step을 한 decoding token으로 simulate하므로
\[
\Pclass\subseteq \CoT(\poly(n)).
\]

Upper bound 쪽에서는 polynomial number of forward passes 각각이 polynomial-time에 simulate 가능하므로
\[
\CoT(\poly(n))\subseteq\Pclass.
\]

따라서 해당 architectural assumptions 아래
\[
\boxed{\CoT(\poly(n))=\Pclass.}
\]

이것은 ``LLM이 P-complete problem을 실제로 학습한다''는 theorem이 아니라,
allowed computation class의 exact characterization이다.

\begin{presentbox}
``CoT의 계산이론적 의미는 intermediate text의 영어 의미가 아닙니다. 중요한 건 token 하나를 쓸 때마다 Transformer를 다시 한 번 실행한다는 사실입니다. 그래서 scratchpad 길이가 serial time budget이 되고, linear CoT부터 automaton을 직접 simulation할 수 있으며 polynomial CoT에서는 P까지 올라갑니다.''
\end{presentbox}

% =====================================================================
\chapter{CoT 길이에 따른 계산력의 계층}
% =====================================================================

\section{\texorpdfstring{$O(1)$: intermediate generation이 없는 경우}{O(1): intermediate generation이 없는 경우}}

Log-precision fixed-depth Transformer 한 번의 출력은
\[
\CoT(O(1))\subseteq\TC^0
\]
관점에서 이해할 수 있다.

따라서 standard separation assumption 아래 automaton simulation처럼 $\NC^1$-complete한 sequential composition 문제를 모두 처리할 수 없다.

\section{\texorpdfstring{$O(\log n)$: 조금의 scratchpad는 얼마나 도움이 되는가?}{O(log n): 조금의 scratchpad는 얼마나 도움이 되는가?}}

Space upper bound
\[
\CoT(t(n))\subseteq\SPACE(t(n)+\log n)
\]
에
\[
t(n)=O(\log n)
\]
을 대입하면
\[
\CoT(O(\log n))\subseteq\SPACE(O(\log n))=\Lclass.
\]

즉 logarithmic CoT는 $\TC^0$보다 약간 더 많은 sequentiality를 줄 수 있지만,
여전히 logspace 범위 안에 있다.

\section{\texorpdfstring{$\Theta(n)$: 모든 regular language}{Theta(n): 모든 regular language}}

Theorem 1의 construction으로
\[
\Reg\subseteq\CoT(\Theta(n)).
\]

동시에 space upper bound는
\[
\CoT(\Theta(n))\subseteq\SPACE(O(n)).
\]

Linear bounded space로 recognize되는 language가 context-sensitive language class와 연결되므로
\[
\CoT(\Theta(n))\subseteq\CSL
\]
라는 upper-bound picture를 얻는다.

\section{\texorpdfstring{$\poly(n)$: 정확히 polynomial time}{poly(n): 정확히 polynomial time}}

앞 절에서 본 것처럼
\[
\CoT(\poly(n))=\Pclass
\]
가 된다.

이를 한 줄 표로 정리하면 다음과 같다.

\begin{center}
\small
\begin{tabularx}{\textwidth}{>{\bfseries\raggedright\arraybackslash}p{0.18\textwidth}Y Y}
\toprule
Intermediate steps & 알려진 계산력 그림 & 해석 \\
\midrule
$O(1)$ & $\subseteq$ uniform $\TC^0$ & highly parallel one-pass computation \\
$O(\log n)$ & $\subseteq \Lclass$ & logspace까지의 작은 sequential workspace \\
$\Theta(n)$ & $\Reg\subseteq\CoT(n)\subseteq\CSL$ & DFA simulation 가능 \\
$\poly(n)$ & $=\Pclass$ (해당 construction assumptions) & polynomial-time TM simulation \\
\bottomrule
\end{tabularx}
\end{center}

\begin{keybox}
CoT의 이론적 효과는 ``reasoning style''이 아니라 \textbf{adaptive sequential depth}를 추가한다는 것이다. 어려운 문제에 더 많은 token을 쓰면 더 많은 serial computation을 배정할 수 있다.
\end{keybox}

\section{왜 graph connectivity 예제가 자주 나오는가?}

Directed graph reachability는 대표적인 $\NL$-complete problem이다.
One-pass $\TC^0$ upper bound와는 큰 gap이 있지만, 충분한 intermediate step을 쓰면 DFS/BFS류 sequential algorithm을 simulate할 수 있다.

ICLR 2024 paper는 random-access가 없는 Turing machine simulation overhead를 고려해 directed connectivity에 quadratic step upper bound를 예시로 든다.
핵심은 구체적인 $n^2$ 상수보다
\[
\text{scratchpad가 graph traversal의 serial state를 저장할 수 있다}
\]
는 점이다.

\begin{presentbox}
``CoT length를 resource axis로 보면 깔끔한 hierarchy가 생깁니다. 상수 step은 $\TC^0$, 로그 step은 logspace upper bound, 선형 step부터 모든 finite automaton을 simulation하고, polynomial step에서는 polynomial-time Turing machine과 맞아집니다. 이게 CoT를 단순 prompting trick이 아니라 compute allocation으로 볼 수 있는 이론적 근거입니다.''
\end{presentbox}

% =====================================================================
\chapter{Automata shortcut: sequential algorithm을 꼭 순차적으로 풀어야 하나?}
% =====================================================================

\section{Liu et al.: shallow Transformer가 automaton을 shortcut할 수 있다}

앞 절에서는 DFA를 token-by-token simulate하려면 $O(n)$ CoT step이면 충분하다고 했다.
그렇다면 CoT 없이 shallow Transformer는 automaton 계산을 전혀 못하는가?

Liu et al. (ICLR 2023)은 더 미묘한 답을 준다.
Automaton transition들의 algebraic structure를 이용하면 recurrent trajectory를 직접 따라가지 않고도 parallel하게 composition할 수 있다.

모든 automaton에 대해 polynomial-size
\[
O(\log n)
\]
depth Transformer simulator가 존재하며, 특정 algebraic structure를 가진 automata에는 $O(1)$-depth shortcut도 존재할 수 있다.

\section{Binary-tree composition}

입력 각 symbol이 state-transition function
\[
f_i:Q\to Q,\qquad f_i(q)=\delta(q,x_i)
\]
를 정한다고 하자.

원래 계산은
\[
q_n=(f_n\circ f_{n-1}\circ\cdots\circ f_1)(q_0).
\]

함수 합성은 associative하므로
\[
(f_4\circ f_3)\circ(f_2\circ f_1)
\]
같은 balanced tree로 계산할 수 있다.

Level 1에서 길이-2 block의 composition,
level 2에서 길이-4 block,
\dots,
level $\lceil\log_2 n\rceil$에서 전체 composition을 얻는다.

따라서 recurrent time $n$을
\[
\boxed{O(\log n)\text{ parallel depth}}
\]
로 줄일 수 있다는 직관이 나온다.

\section{\texorpdfstring{이 결과와 $\TC^0$ upper bound는 충돌하지 않는다}{이 결과와 TC0 upper bound는 충돌하지 않는다}}

$O(\log n)$ depth는 constant depth가 아니다.
Input length에 따라 Transformer layer 수 자체를 키우는 family를 허용하므로 one-pass fixed-depth $\TC^0$ theorem과 setting이 다르다.

또 일부 automaton은 constant-depth shortcut을 갖지만, \textbf{모든} automaton에 대해 그런 shortcut이 존재하는 것은 아니다.

\section{Learning과 existence를 다시 분리한다}

Liu et al.의 실험은 synthetic automata에서 training이 shortcut solution을 찾을 수 있음을 보여 주지만,
이런 solution이 length extrapolation에서 brittle할 수 있음을 함께 지적한다.

따라서
\[
\text{existence of shallow circuit}
\not\Rightarrow
\text{robust algorithmic generalization}.
\]

이 구분은 Week 4의 UAT와도 같은 패턴이다.
표현 가능성 theorem과 실제 optimization/generalization theorem은 다르다.

\begin{presentbox}
``Sequential algorithm이라고 해서 반드시 $n$ layer가 필요한 것은 아닙니다. Automaton transition은 associative composition이므로 tree reduction으로 $O(\log n)$ depth shortcut을 만들 수 있습니다. 다만 그건 layer 수가 input length와 함께 증가하는 construction이고, 학습된 shortcut이 길이 밖에서도 robust하다는 보장은 별도 문제입니다.''
\end{presentbox}

% =====================================================================
\chapter{Formal language hierarchy로 결과를 읽는 법}
% =====================================================================

\section{대표 언어 세 개}

\subsection{Parity}
\[
\mathcal L_{\mathrm{parity}}
=
\left\{x\in\bits^*: \sum_i x_i\equiv 1\pmod 2\right\}.
\]
Regular language지만 $\AC^0$ 밖에 있다. 따라서 ``regular''이 곧 ``constant-depth AND/OR circuit로 쉽다''는 뜻은 아니다.

\subsection{Dyck-1}
한 종류의 괄호에 대해 prefix balance
\[
b_k=\sum_{i=1}^{k}
\begin{cases}
+1,&x_i=(,\\
-1,&x_i=),
\end{cases}
\]
가
\[
\min_{k\le n} b_k\ge 0,
\qquad
b_n=0
\]
을 만족하는 language다.
Unbounded nesting 때문에 context-free structure의 대표 예로 쓰인다.

\subsection{Automaton evaluation}
Automaton $A$와 string $x$가 주어졌을 때
\[
\delta^*(q_0,x)\in F?
\]
를 판정한다.
Automaton이 input의 일부로 주어지는 uniform simulation problem은 단순히 하나의 고정 regular language를 판정하는 것보다 훨씬 강한 계산 문제다.

\section{``Regular language를 못한다''는 표현을 조심해야 하는 이유}

어떤 fixed Transformer가 특정 regular language를 recognize할 수 있는 것과
\[
\text{모든 regular language를 하나의 architecture class가 recognize할 수 있는가}
\]
는 다르다.

예를 들어 parity 자체도 regular지만 hard-attention $\AC^0$ upper bound 아래에서는 불가능하다.
반면 ``마지막 symbol이 1인가?'' 같은 regular language는 trivial한 local rule이다.

따라서 발표에서는
\begin{itemize}
  \item ``Transformer는 regular language를 못한다''가 아니라,
  \item ``해당 fixed-resource Transformer class는 모든 regular language를 cover하지 못한다''
\end{itemize}
라고 말하는 것이 안전하다.

\section{Turing completeness와 formal-language upper bound를 동시에 기억하는 법}

다음 세 문장을 따로 기억하면 된다.

\begin{enumerate}
  \item \textbf{Fixed-depth one-pass:} circuit class로 upper-bound될 수 있다.
  \item \textbf{Depth grows with $n$:} tree-style parallel algorithms이 가능해진다.
  \item \textbf{Autoregressive steps grow with $n$:} serial computation을 직접 펼칠 수 있다.
\end{enumerate}

같은 Transformer block을 쓰더라도 실행 graph가 달라지면 complexity class도 달라진다.

\begin{presentbox}
``Formal-language 결과를 읽을 때 language 이름보다 quantifier를 보세요. 특정 regular language를 풀 수 있느냐, 모든 regular language를 cover하느냐, automaton 자체가 input으로 주어지느냐가 서로 다른 문제입니다. 이걸 구분하면 parity, DFA simulation, Turing completeness 결과가 한 그림 안에 들어옵니다.''
\end{presentbox}

% =====================================================================
\chapter{Arithmetic와 algorithmic reasoning을 어떻게 해석할까?}
% =====================================================================

\section{\texorpdfstring{$\TC^0$ 안에도 강한 arithmetic이 있다}{TC0 안에도 강한 arithmetic이 있다}}

Threshold circuits는 단순 boolean pattern matching보다 강하다.
정수 addition, multiplication, iterated arithmetic의 상당 부분이 uniform $\TC^0$에 들어간다는 고전 결과가 있다.

따라서 log-precision Transformer가 한 pass에서
\[
\text{weighted sums, normalization, counting-like arithmetic}
\]
을 수행하는 것 자체는 $\TC^0$ upper bound와 전혀 모순되지 않는다.

\section{어려움은 arithmetic operation 자체보다 sequential dependency일 수 있다}

예를 들어 $n$개의 수를 더하는 것은 highly parallelizable하다.
반면 state가 매 step 바뀌며 다음 operation이 이전 결과에 의존하는 algorithm은
\[
s_{t+1}=F(s_t,x_t)
\]
꼴의 serial dependency를 가진다.

Transformer가 algorithmic reasoning에서 어려움을 겪는 이유를 단순히 ``산수를 못해서''라고 보는 것보다
\[
\boxed{\text{필요한 adaptive computational depth가 충분한가?}}
\]
로 보는 관점이 더 이론과 잘 맞는다.

\section{CoT가 제공하는 것은 externalized recurrence}

Autoregressive decoder에서는
\[
y_{t+1}=f(x,y_{1:t})
\]
이므로 prefix가 recurrent state 역할을 한다.

일반 RNN의
\[
h_{t+1}=G(h_t,x_t)
\]
와 비교하면, Transformer는 state를 fixed-size hidden vector에 압축하지 않고
\[
(y_1,\dots,y_t)
\]
라는 growing external memory에 저장한다.

Attention은 이 history에서 필요한 record를 random-access에 가깝게 retrieve하고,
FFN은 local transition을 계산한다.

\begin{suppbox}
이 관점에서 CoT의 계산력 증가는 ``문장을 자연어로 설명해서''가 아니라 다음 구조 때문이다.
\[
\begin{aligned}
\text{write token}
&\rightarrow \text{append memory}
\rightarrow \text{next forward pass}\\
&\rightarrow \text{read old memory}
\rightarrow \text{write next token}.
\end{aligned}
\]
즉 external-memory recurrent machine이 만들어진다.
\end{suppbox}

\section{그러면 내부 hidden reasoning은 의미가 없는가?}

그렇지 않다. 한 forward pass 안에서도 각 layer는 $L$단계의 computation을 한다.
다만 $L$이 input length와 무관한 상수라면 asymptotic circuit depth는 고정된다.

따라서 실제 reasoning capacity를 크게 세 축으로 나눌 수 있다.
\[
\text{within-pass depth }L,
\qquad
\text{between-pass steps }t(n),
\qquad
\text{memory/precision }p(n).
\]

\begin{presentbox}
``$\TC^0$는 arithmetic을 못하는 class가 아닙니다. 오히려 병렬 arithmetic은 상당히 잘합니다. 문제는 이전 계산 결과를 다음 step에서 다시 써야 하는 긴 dependency chain입니다. CoT는 output prefix를 external state로 만들어 이 sequential depth를 모델 밖으로 늘려 줍니다.''
\end{presentbox}

% =====================================================================
\chapter{결과를 과장하지 않기 위한 수학적 체크리스트}
% =====================================================================

\section{\texorpdfstring{Turing complete $\neq$ efficient}{Turing complete != efficient}}

Turing completeness는 존재성 결과다.
어떤 computation을 simulate할 수 있다는 것과
\[
\text{time},\text{space},\text{precision},\text{learnability}
\]
가 실용적이라는 것은 별개다.

특히 arbitrary-precision rational을 사용하는 proof는 finite-precision GPU implementation과 직접 동일시하면 안 된다.

\section{\texorpdfstring{$\TC^0$ upper bound $\neq$ empirical impossibility}{TC0 upper bound != empirical impossibility}}

Asymptotic exact recognition theorem은
\[
\forall n,\ \forall x\in\Sigma^n
\]
에 대한 statement다.

실제 benchmark는 finite $n$, finite distribution, approximate accuracy를 본다.
따라서 complexity lower bound가 있어도 training distribution 안에서는 shortcut/heuristic으로 높은 accuracy가 가능하다.

\section{\texorpdfstring{CoT power $\neq$ faithful reasoning}{CoT power != faithful reasoning}}

Computational expressivity theorem은 generated token이 인간이 해석하는 reasoning trace와 일치한다는 것을 보장하지 않는다.
중간 token은 계산 resource일 수 있지만, 그 자연어 내용의 faithful explanation 여부는 별도 interpretability 문제다.

\section{Architecture assumption을 반드시 붙인다}

다음 문장들은 각각 qualifier가 필요하다.

\begin{center}
\small
\begin{tabularx}{\textwidth}{>{\bfseries\raggedright\arraybackslash}p{0.30\textwidth}Y}
\toprule
짧게 말하면 위험한 문장 & 정확한 버전 \\
\midrule
Transformer는 $\TC^0$다. & log-precision, fixed-depth one-pass Transformer classifier에 대한 upper bound다. \\
Transformer는 Turing complete다. & arbitrary precision, positional encoding, hard-attention, unbounded decoding을 허용한 Pérez et al. 모델에서 성립한다. \\
CoT가 있으면 P를 푼다. & polynomially many decoding steps와 특정 architectural construction 아래 recognition class가 $\Pclass$와 일치한다. \\
Transformer는 regular language를 못한다. & 특정 fixed-resource classes는 모든 regular language를 cover하지 못한다. \\
Global attention이면 sequential task도 된다. & global communication과 computational depth는 별개다. \\
\bottomrule
\end{tabularx}
\end{center}

\section{세미나에서 가장 중요한 conceptual distinction}

\begin{align*}
\text{Expressivity}
&:\quad \exists\theta\text{가 target function을 표현하는가?}\\
\text{Computational complexity}
&:\quad \text{input length에 따라 필요한 depth/time/space는?}\\
\text{Learnability}
&:\quad \text{optimization이 그 }\theta\text{를 찾는가?}\\
\text{Generalization}
&:\quad \text{더 긴 입력이나 새 distribution에도 유지되는가?}
\end{align*}

Week 4의 UAT는 첫 번째 질문에 가까웠고, 이번 주는 두 번째 질문을 다룬다.

\begin{presentbox}
``이번 주 theorem은 실제 GPT가 어떤 문제를 맞히는지 직접 예측하는 실험 법칙이 아닙니다. 무엇을 resource로 세느냐에 따라 architecture가 어느 complexity class까지 갈 수 있는지를 분류하는 결과입니다. 그래서 theorem을 인용할 때 attention, precision, depth, decoding steps를 같이 말하는 게 가장 중요합니다.''
\end{presentbox}

% =====================================================================
\chapter{발표 준비용 압축 정리}
% =====================================================================

\section{칠판에 반드시 남길 열 개 식}

\textbf{1. DFA recurrence}
\[
q_i=\delta(q_{i-1},x_i).
\]

\textbf{2. Circuit hierarchy}
\[
\AC^0\subsetneq\TC^0\subseteq\NC^1\subseteq\Lclass\subseteq\NL\subseteq\Pclass.
\]

\textbf{3. Log-precision Transformer upper bound}
\[
\text{fixed depth}+O(\log n)\text{-bit precision}
\subseteq
\text{uniform }\TC^0.
\]

\textbf{4. TM transition}
\[
\delta(q_t,\tau_t(h_t))
=(q_{t+1},s_t^{\mathrm{write}},d_t).
\]

\textbf{5. Pérez positional codes}
\[
\operatorname{pos}(n)\text{ uses }n,\frac1n,\frac1{n^2}.
\]

\textbf{6. CoT iteration}
\[
y_{t+1}=f(x,y_{1:t}).
\]

\textbf{7. CoT class}
\[
\TIME(t(n))
\subseteq
\CoT(t(n))
\subseteq
\SPACE(t(n)+\log n).
\]

\textbf{8. Layer-norm hash}
\[
\phi_x=
\frac{\langle x,1,-x,-1\rangle}{\sqrt{2x^2+2}}.
\]

\textbf{9. Equality retrieval}
\[
\inner{\phi_q}{\phi_k}=1\iff q=k.
\]

\textbf{10. Polynomial scratchpad}
\[
\CoT(\poly(n))=\Pclass
\quad\text{(paper의 construction assumptions 아래).}
\]

\section{발표 순서 추천}

\begin{enumerate}
  \item Week 4와 연결해 expressivity와 computational complexity를 분리한다.
  \item formal language와 $\AC^0/\TC^0$를 최소한으로 정의한다.
  \item hard-attention / Hahn 결과로 fixed-depth limitation intuition을 준다.
  \item Merrill--Sabharwal의 log-precision $\TC^0$ theorem을 설명한다.
  \item Pérez Turing-complete theorem을 보여 주고 일부러 ``모순처럼'' 제시한다.
  \item precision과 decoding time이 다르므로 모순이 아님을 표로 해소한다.
  \item CoT를 serial forward-pass budget으로 정의한다.
  \item layer-norm hash를 직접 유도한다.
  \item $O(1)$, $O(\log n)$, $O(n)$, $\poly(n)$ hierarchy를 설명한다.
  \item 마지막에 learnability/generalization과 분리한다.
\end{enumerate}

\section{자주 나오는 질문}

\textbf{Q1. Transformer가 Turing complete라면 왜 parity 같은 단순 문제에 limitation theorem이 있나?}

서로 다른 모델이다. Turing-complete proof는 arbitrary precision과 unbounded decoder steps를 쓴다. Fixed-depth/limited-precision language recognizer theorem은 그런 resource를 제한한다.

\textbf{Q2. $\TC^0$면 덧셈이나 곱셈도 못하는 약한 class인가?}

아니다. Threshold circuits는 강한 parallel arithmetic을 포함한다. 한계는 주로 긴 serial dependency를 요구하는 computation에서 나타난다.

\textbf{Q3. CoT가 길수록 항상 실제 accuracy가 좋아진다는 theorem인가?}

아니다. 더 긴 CoT가 더 큰 computation class를 \textbf{표현할 수 있게} 한다는 결과다. Optimization, error accumulation, faithfulness는 별개다.

\textbf{Q4. Linear CoT가 모든 regular language를 푼다면 one-pass Transformer는 regular language를 하나도 못 푸나?}

아니다. 쉬운 regular language는 one-pass에서도 가능하다. Linear CoT result는 \textbf{모든} regular language에 대한 coverage를 준다.

\textbf{Q5. Turing-complete proof에 positional encoding이 왜 필요한가?}

Unbounded computation history에서 exact time/address 정보를 구분해야 하기 때문이다. Pérez construction은 rational positional codes로 이것을 구현한다.

\textbf{Q6. Layer-norm이 계산력을 늘리는 게 이상하지 않나?}

일반적 의미로 layer-norm이 항상 계산력을 늘린다는 주장이 아니다. Merrill--Sabharwal construction에서는 scale invariance
\[
\phi(x/i,1/i)=\phi(x,1)
\]
를 이용해 finite precision에서 address equality를 구현하는 specific gadget로 사용한다.

\textbf{Q7. Automaton은 finite memory인데 왜 one-pass Transformer가 어려워할 수 있나?}

Memory size와 parallel computational depth가 다르기 때문이다. State는 작아도 $n$개의 transition composition을 계산해야 한다.

\section{흔한 수학적 실수 체크리스트}

\begin{itemize}
  \item arbitrary precision theorem과 log-precision theorem을 같은 Transformer model로 취급하지 않았는가?
  \item one forward pass와 autoregressive $t(n)$ passes를 구분했는가?
  \item $\AC^0$와 $\TC^0$를 혼동하지 않았는가?
  \item $\TC^0\neq\NC^1$ 같은 separation을 proven fact처럼 말하지 않았는가?
  \item ``모든 regular languages''와 ``어떤 regular language''를 구분했는가?
  \item Turing completeness를 efficiency/learnability로 해석하지 않았는가?
  \item CoT($\poly(n)$) $=\Pclass$의 architectural assumptions를 생략하지 않았는가?
  \item layer-norm hash의 equality condition에서 Cauchy--Schwarz equality를 정확히 썼는가?
  \item $O(\log n)$ precision과 $O(1)$ fixed precision을 섞지 않았는가?
  \item formal-language exact recognition을 finite-distribution accuracy와 동일시하지 않았는가?
\end{itemize}

\section{30초 결론}

Fixed-depth Transformer 한 번의 forward pass는 매우 parallel한 계산이다. Log-precision setting에서는 그 계산을 uniform $\TC^0$ threshold circuit로 simulate할 수 있어 긴 sequential computation에 asymptotic limitation이 생긴다. 반면 Pérez et al.의 arbitrary-precision hard-attention Transformer는 autoregressive history를 memory로 사용하면서 Turing machine의 step을 계속 simulate할 수 있어 Turing complete하다. 둘은 precision과 sequential-time assumptions가 다르므로 모순이 아니다. Merrill--Sabharwal은 이 사이를 chain-of-thought length로 정량화해, intermediate generation이 $t(n)$번의 serial forward pass를 제공하며 linear steps에서 모든 regular language, polynomial steps에서 적절한 architecture 가정 아래 정확히 $\Pclass$의 계산력에 도달함을 보였다.

\begin{presentbox}
``한 문장으로 정리하면, Transformer의 핵심 tradeoff는 parallelism과 sequential compute입니다. 한 번의 고정-depth pass는 매우 parallel해서 $\TC^0$로 제한될 수 있지만, output token을 다시 input으로 넣는 순간 token 수만큼 계산 depth를 외부로 늘릴 수 있습니다. CoT가 computational resource라는 말은 정확히 이 뜻입니다.''
\end{presentbox}

% =====================================================================
\chapter*{참고문헌 및 원문에서 확인할 위치}
\addcontentsline{toc}{chapter}{참고문헌 및 원문에서 확인할 위치}
% =====================================================================

\textbf{Jorge Pérez, Pablo Barceló, Javier Marinkovic.}
``Attention is Turing Complete.'' Journal of Machine Learning Research 22, 2021.

\begin{itemize}
  \item Sec. 2--3: Turing completeness definition, Transformer formalization, proportion invariance.
  \item Theorem 6: positional encodings을 가진 Transformer의 Turing completeness.
  \item Theorem 6 statement: positional code의 nonconstant terms $n,1/n,1/n^2$, one encoder + three decoder layers.
  \item Sec. 4: Turing-machine simulation construction.
  \item Sec. 5: hard attention과 arbitrary precision이 proof에서 수행하는 역할; fixed precision limitation.
\end{itemize}
\url{https://jmlr.org/papers/v22/20-302.html}

\vspace{0.8em}
\textbf{William Merrill, Ashish Sabharwal.}
``The Parallelism Tradeoff: Limitations of Log-Precision Transformers.'' TACL 11, 2023.

\begin{itemize}
  \item Main theorem: $O(\log n)$ precision Transformer를 logspace-uniform constant-depth threshold circuit로 simulate.
  \item Sec. 2: $\TC^0$ upper bound의 computational interpretation과 parallelism tradeoff.
  \item Proof sections: float/arithmetic primitives와 Transformer block의 threshold-circuit simulation.
\end{itemize}
\url{https://aclanthology.org/2023.tacl-1.31/}

\vspace{0.8em}
\textbf{William Merrill, Ashish Sabharwal.}
``The Expressive Power of Transformers with Chain of Thought.'' ICLR 2024.

\begin{itemize}
  \item Eq. (1): $\TIME(t(n))\subseteq\CoT(t(n))\subseteq\SPACE(t(n)+\log n)$와 time upper bound.
  \item Theorem 1: $|x|+1$ decoding steps로 모든 regular language recognition.
  \item Theorem 2: one decoding step per Turing-machine step simulation.
  \item Lemma 1--2: layer-norm hash와 equality check.
  \item Theorem 3--4: time/space upper bounds.
  \item Sec. 4--5: logarithmic, linear, polynomial CoT의 complexity interpretation.
\end{itemize}
\url{https://proceedings.iclr.cc/paper_files/paper/2024/hash/1f59721c106ea80f613299039112f651-Abstract-Conference.html}

\vspace{0.8em}
\textbf{Michael Hahn.}
``Theoretical Limitations of Self-Attention in Neural Sequence Models.'' TACL 8, 2020.

\begin{itemize}
  \item Parity와 Dyck를 이용한 fixed-resource self-attention limitation.
  \item Input length에 따라 layer/head resource를 키우지 않을 때의 formal-language limitations.
\end{itemize}
\url{https://aclanthology.org/2020.tacl-1.11/}

\vspace{0.8em}
\textbf{Yiding Hao, Dana Angluin, Robert Frank.}
``Formal Language Recognition by Hard Attention Transformers: Perspectives from Circuit Complexity.'' TACL 10, 2022.

\begin{itemize}
  \item UHAT/GUHAT의 $\AC^0$ upper bound.
  \item AHAT가 MAJORITY와 DYCK-1을 recognize할 수 있는 separation.
\end{itemize}
\url{https://aclanthology.org/2022.tacl-1.46/}

\vspace{0.8em}
\textbf{Bingbin Liu, Jordan T. Ash, Surbhi Goel, Akshay Krishnamurthy, Cyril Zhang.}
``Transformers Learn Shortcuts to Automata.'' ICLR 2023.

\begin{itemize}
  \item 모든 finite automaton에 대한 polynomial-size $O(\log n)$-depth shortcut construction.
  \item algebraic structure에 따른 shallower construction과 empirical shortcut learning.
  \item length generalization에서 shortcut brittleness에 대한 실험적 분석.
\end{itemize}
\url{https://openreview.net/forum?id=De4FYqjFueZ}

\begin{warningbox}[--- 원문과 이 노트의 경계]
Circuit-class 정의, DFA transition을 balanced composition으로 쓰는 $O(\log n)$ 직관, Pérez construction을 ``event log + latest write lookup''으로 설명한 부분, CoT를 externalized recurrence로 해석한 부분은 원 theorem을 이해하기 위한 보충 전개다. Pérez et al.과 Merrill--Sabharwal의 full simulation proof는 많은 coordinate-level gadget을 포함하므로 이 노트는 theorem statement와 proof mechanism을 보존하되 모든 coordinate construction을 복제하지 않았다. 또한 $\TC^0\neq\NC^1$, $\Lclass\neq\Pclass$ 같은 separation은 표준 conjecture이지 증명된 사실로 다루지 않았다.
\end{warningbox}

\end{document}
